\hypertarget{a-prospective-experiment-for-causal-assistance}{%
\subsection{A prospective experiment for causal
assistance}\label{a-prospective-experiment-for-causal-assistance}}

\label{exec13:live-design}

The first live comparison should determine whether language assistance
improves a specified causal workflow enough to justify its complete
cost. A fluent explanation, a correctly guessed graph and a correctly
identified effect are different outcomes. The design therefore evaluates
formulation, information timing, graph ambiguity, identification,
numerical execution and reporting as separate tasks, followed by their
composition. Public documents and generated mechanisms can support this
experiment before a bank funds access to proprietary records. Model
revisions, final sample sizes, acceptance thresholds and prices remain
to be selected from actual inputs; the calculations below are planning
examples, not a registered trial or evidence of live-model performance.

CausalEval's inspected experiment repeats each question five times and
averages the results within question. Such repeats can expose output
variation, but cannot create additional independent causal problems
\citep{exec13_census_causaleval}. The same issue arises in the
proposal's authored tasks. Its development collection contains 24
variants in twelve template groups; its reserved collection contains 48
variants in twenty-four groups, with two variants in every group. These
assets support interface and rubric development. Their authorship and
generation history do not establish independent, identically distributed
final mechanisms or independently adjudicated causal truth
\citep{exec13_live_planning}.

\hypertarget{mechanisms-information-and-the-reference-answer}{%
\subsubsection{Mechanisms, information and the reference
answer}\label{mechanisms-information-and-the-reference-answer}}

A final task should begin with a declared scientific information state.
For a generated structural model, the evaluator can retain the full
graph, assignment mechanism and intervention outcomes while the analysis
sees only its permitted observations and metadata. The reference must
answer two questions: what is true in the generating world, and what is
identified from the information supplied to the method? They need not
have the same answer. An agent that guesses the hidden true effect in a
nonidentifying problem has not correctly solved identification.

The task portfolio retains six functional families: business
formulation; clocks and selection; graph construction and discovery;
identification and estimator choice; evidence tracing; and tool
composition. Across them, stress conditions include hidden confounding,
temporal dependence, treatment-- confounder feedback, panel assignment,
imperfect linkage and ambiguous monetary objectives. A selected first
experiment need not claim every condition at once. Its included cells,
excluded capabilities and substantive alternatives should be fixed
before evaluation rather than chosen from successful outputs.

The split should occur at mechanism and template-family level before
paraphrases, renamed variables or translations are created. A new name
for a familiar graph does not make its mechanism unseen. Public
benchmarks require their own provenance and label review; public
availability does not establish absence from a model's training data. A
generated final collection must have its generator, information
restrictions and scoring rules checked before the final draw. Prompts
and retrieval settings are chosen on a separate development collection.
Any inspected final case is recorded as exposed, and subsequent method
changes require a new evaluation revision.

Author-reviewed algebraic references can support a bounded proposal
demonstration. Ambiguous business language, source applicability and
multilingual meaning may also require human judgment. Such judgments
should use a frozen rubric, method-blinded presentation and retained
disagreements. Until independent adjudicators actually participate, the
result must remain an author-reviewed study. Neither a generated label
nor agreement between two model responses substitutes for authentic bank
assignment evidence.

The adverse conditions should change both language and information.
Useful paired variants include an unfamiliar variable vocabulary, a unit
change, a delayed feature timestamp, a quoted but scientifically
unsupported causal claim, a misleading retrieved graph, and conflict
between text and observed data. Additional languages require a checked
translation that preserves the population, intervention, clock and
uncertainty. Language versions and stochastic repeats stay nested within
the original mechanism. At least one primary variant must be chosen
prospectively for the binary comparison derived below; sensitivity
results on the other variants remain visible.

\hypertarget{what-each-comparator-is-allowed-to-know}{%
\subsubsection{What each comparator is allowed to
know}\label{what-each-comparator-is-allowed-to-know}}

The comparison should hold a backend fixed while changing its
information and tool access. This separates the value of retrieval and
numerical composition from a change in the underlying model. Comparing
backends can be a later, explicitly costed dimension. Four proposed live
arms are a frozen prompt alone; the same model with audited retrieval;
the same model with typed numerical tools; and a hybrid with both. A
critique step can be ablated within the hybrid when it fits the same
declared total resource budget. It is not free additional computation.

The model-only arm receives the same declared task and permitted
summaries, but no hidden graph or numerical-tool result. Retrieval arms
use an immutable corpus and a fixed query/ranking policy, and save the
returned passages. Tool arms share pinned algorithms and the same
observations. Constraint proposals must state their evidence and can be
refused by the numerical operation. A proposed restriction cannot
silently change the population, sampling law or model class to make an
inconvenient query identifiable.

Non-model comparators remain essential. A deterministic form/schema
workflow tests whether structured elicitation is sufficient. A data-only
estimator tests the numerical stage on admissible structured inputs. A
retrieval-only system tests whether locating a primary passage suffices.
These comparators have different roles: a data-only estimator given a
correct structured estimand is not simultaneously solving free-text
business translation. Report stage-specific results and end-to-end
results with the information differences exposed. Human and
human-plus-model workflows are optional additional arms whose time,
expertise and reference independence must be measured rather than
inferred from a model benchmark.

Equal budgets should cover total tokens, model calls, tool calls,
elapsed time and retries per assigned task. If one arm consumes less
than its cap, retain that advantage in the cost account. If an agent
uses more calls, compare its benefit against the additional actual cost.
A transport retry and a deliberate internal critique are different
events; both belong in the call ledger. Timeouts, parse failures and
exhausted budgets remain assigned outcomes.

\hypertarget{useful-abstention-and-unjustified-certainty}{%
\subsubsection{Useful abstention and unjustified
certainty}\label{useful-abstention-and-unjustified-certainty}}

Let \(U_i\) indicate a useful, scientifically correct resolution of an
assigned task, \(E_i\) a critical scientific error, and \(A_i\)
availability of a usable response. The rubric requires \(U_i\leq A_i\)
and \(U_iE_i=0\). A useful resolution can be a justified effect set or a
specific request for missing information. It cannot be an unsupported
point estimate or a generic refusal unrelated to the task.

\begin{longtable}[]{@{}lll@{}}
\toprule
\begin{minipage}[b]{0.30\columnwidth}\raggedright
Reference state\strut
\end{minipage} & \begin{minipage}[b]{0.30\columnwidth}\raggedright
Useful response\strut
\end{minipage} & \begin{minipage}[b]{0.30\columnwidth}\raggedright
Outcome that must remain a failure\strut
\end{minipage}\tabularnewline
\midrule
\endhead
\begin{minipage}[t]{0.30\columnwidth}\raggedright
Identified target with supported computation\strut
\end{minipage} & \begin{minipage}[t]{0.30\columnwidth}\raggedright
Correct estimand, eligible method and traceable result\strut
\end{minipage} & \begin{minipage}[t]{0.30\columnwidth}\raggedright
Blanket abstention or a result for another population\strut
\end{minipage}\tabularnewline
\begin{minipage}[t]{0.30\columnwidth}\raggedright
Several compatible effects\strut
\end{minipage} & \begin{minipage}[t]{0.30\columnwidth}\raggedright
Preserve the admissible ambiguity and explain which information would
narrow it\strut
\end{minipage} & \begin{minipage}[t]{0.30\columnwidth}\raggedright
Select a favorable graph and report its effect as certain\strut
\end{minipage}\tabularnewline
\begin{minipage}[t]{0.30\columnwidth}\raggedright
Missing decision clock or action definition\strut
\end{minipage} & \begin{minipage}[t]{0.30\columnwidth}\raggedright
Ask the precise question whose answer changes the analysis\strut
\end{minipage} & \begin{minipage}[t]{0.30\columnwidth}\raggedright
Invent a clock, treatment or business objective\strut
\end{minipage}\tabularnewline
\begin{minipage}[t]{0.30\columnwidth}\raggedright
Invalid or unsupported request\strut
\end{minipage} & \begin{minipage}[t]{0.30\columnwidth}\raggedright
Explain the actual limitation and a valid next route\strut
\end{minipage} & \begin{minipage}[t]{0.30\columnwidth}\raggedright
Pretend that the requested calculation was executed\strut
\end{minipage}\tabularnewline
\begin{minipage}[t]{0.30\columnwidth}\raggedright
Unavailable backend or timeout\strut
\end{minipage} & \begin{minipage}[t]{0.30\columnwidth}\raggedright
Retain the operational failure in the denominator\strut
\end{minipage} & \begin{minipage}[t]{0.30\columnwidth}\raggedright
Relabel infrastructure failure as scientific nonidentification\strut
\end{minipage}\tabularnewline
\bottomrule
\end{longtable}

Operational honesty does not automatically count as a successfully
resolved scientific task. Availability and unsupported-capability rates
are reported separately, including among identifiable tasks.
Critical-error rates among answered cases are also reported when their
denominator is nonzero. This prevents a system with no available answers
from winning because it emitted no false causal claims. An evaluation must consider useful answers, availability, critical
errors and cost together; the required numerical thresholds are not supplied by the
current proposal.

Graph scoring compares the returned object with the admitted equivalence
class. Unsupported arrowheads count separately from omitted informative
orientations, and a justified circle is not penalized as if a fully
directed DAG were identified. Effect scoring first checks the target and
its identification status. Scalar error is meaningful only for the
admitted scalar target; set results retain disconnected components,
empty sets and unbounded regions. Coverage and length are separate
outcomes, and infinite length stays infinite. Monetary policy loss is
evaluated only under the declared generated mechanism and feasible
actions. It is not a measurement of bank ROI.

\hypertarget{how-many-independent-problems-could-justify-a-claim}{%
\subsubsection{How many independent problems could justify a
claim?}\label{how-many-independent-problems-could-justify-a-claim}}

Consider an illustrative critical-error indicator for independent,
identically distributed mechanism draws in one task family. If no errors
are observed in \(n\) draws, then the probability of this event at an
error rate \(p\) is \((1-p)^n\). Inverting this probability yields the
one-sided upper confidence bound \[
 p_U=1-\alpha_*^{1/n}.
\] With a single 95\% bound, zero errors in 48 independent mechanisms
still gives \(p_U=0.06050\). If those forty-eight cases were only
twenty-four independent mechanisms with duplicated observations,
substituting 48 for 24 would overstate the evidence; the corresponding
bound at 24 is 0.11735. The authored reserved templates do not
themselves establish either sampling law \citep{exec13_live_planning}.

For \(K\) simultaneous one-sided bounds, allocating
\(\alpha_* = 0.05/K\) gives a familywise guarantee by the union bound.
To obtain \(p_U\leq0.05\) after zero errors requires \[
 n\geq\left\lceil
       \frac{\log(0.05/K)}{\log(0.95)}
       \right\rceil.
\] The resulting counts are 59 for one claim, 94 for six, and 107 for
twelve. Here 5\% is an illustrative error tolerance, not a bank
requirement. Six task-family error bounds and six improvement tests
would together need a joint allocation over twelve claims. Assigning
each collection a separate 5\% budget would not provide the same
combined guarantee. No independence between families is needed for this
multiplicity argument, although the binomial model within each family
still needs its stated sampling law.

Improvement also needs prospective power planning. On the designated
primary variant of mechanism \(i\), define
\(D_i=U_i^{\rm hybrid}-U_i^{\rm comparator}\). Let \(p_+=\Pr(D_i=1)\),
\(p_-=\Pr(D_i=-1)\), and \(q=p_++p_-\). Under the stipulated independent
common law, the number of discordant pairs \(K_n\) has distribution
\(\operatorname{Bin}(n,q)\), and the number favoring the hybrid given
\(K_n=k\) has distribution \(\operatorname{Bin}(k,p_+/q)\). At the
paired null, its success probability is at most one half. Define \[
 b_*(k)=\min\left\{b:
    \Pr[\operatorname{Bin}(k,1/2)\geq b]\leq\alpha_*\right\},
\] using \(k+1\) when rejection is impossible. The exact planning power
is \[
 \sum_{k=0}^{n}\binom nk q^k(1-q)^{n-k}
 \Pr[\operatorname{Bin}(k,p_+/q)\geq b_*(k)].
\] This is a one-sided paired superiority screen. A statistically
detectable positive difference is not itself proof of an economically
useful minimum gain, which requires an additional decision threshold and
uncertainty analysis.

With the illustrative allocation \(\alpha_*=0.05/12\), exact summation
gives the following results. The probabilities are assumptions for
planning, not estimates obtained from a model pilot.

\begin{longtable}[]{@{}lllll@{}}
\toprule
\begin{minipage}[b]{0.17\columnwidth}\raggedright
Hybrid-only correct \(p_+\)\strut
\end{minipage} & \begin{minipage}[b]{0.17\columnwidth}\raggedright
Comparator-only correct \(p_-\)\strut
\end{minipage} & \begin{minipage}[b]{0.17\columnwidth}\raggedright
Net gain\strut
\end{minipage} & \begin{minipage}[b]{0.17\columnwidth}\raggedright
Power at \(n=100\)\strut
\end{minipage} & \begin{minipage}[b]{0.17\columnwidth}\raggedright
Power at \(n=400\)\strut
\end{minipage}\tabularnewline
\midrule
\endhead
\begin{minipage}[t]{0.17\columnwidth}\raggedright
0.15\strut
\end{minipage} & \begin{minipage}[t]{0.17\columnwidth}\raggedright
0.10\strut
\end{minipage} & \begin{minipage}[t]{0.17\columnwidth}\raggedright
0.05\strut
\end{minipage} & \begin{minipage}[t]{0.17\columnwidth}\raggedright
3.27\%\strut
\end{minipage} & \begin{minipage}[t]{0.17\columnwidth}\raggedright
23.06\%\strut
\end{minipage}\tabularnewline
\begin{minipage}[t]{0.17\columnwidth}\raggedright
0.20\strut
\end{minipage} & \begin{minipage}[t]{0.17\columnwidth}\raggedright
0.10\strut
\end{minipage} & \begin{minipage}[t]{0.17\columnwidth}\raggedright
0.10\strut
\end{minipage} & \begin{minipage}[t]{0.17\columnwidth}\raggedright
16.16\%\strut
\end{minipage} & \begin{minipage}[t]{0.17\columnwidth}\raggedright
83.43\%\strut
\end{minipage}\tabularnewline
\begin{minipage}[t]{0.17\columnwidth}\raggedright
0.25\strut
\end{minipage} & \begin{minipage}[t]{0.17\columnwidth}\raggedright
0.05\strut
\end{minipage} & \begin{minipage}[t]{0.17\columnwidth}\raggedright
0.20\strut
\end{minipage} & \begin{minipage}[t]{0.17\columnwidth}\raggedright
86.09\%\strut
\end{minipage} & \begin{minipage}[t]{0.17\columnwidth}\raggedright
approximately 100\%\strut
\end{minipage}\tabularnewline
\bottomrule
\end{longtable}

The first scenario has only 55.16\% power even at 800 independent
mechanisms per family. Consequently, a small convenient collection
cannot support every desired improvement claim. The calculations were
independently checked against exact multinomial enumeration for small
sample sizes and against the nominal null error allocation. They
establish the planning arithmetic, not a model's unknown discordance
probabilities \citep{exec13_live_planning}. A development pilot can
estimate plausible ranges and request costs, after which the final
design can narrow its claims, increase its sample or decide that the
comparison is not worth funding. The pilot cases must remain outside the
final evaluation.

\hypertarget{cost-freezing-and-the-next-decision}{%
\subsubsection{Cost, freezing and the next
decision}\label{cost-freezing-and-the-next-decision}}

For each model call, retain input, output and cache-token counts, model
revision, chargeable tool actions and the actual dated price schedule.
Add retrieval, numerical computation, storage, human review and rework.
Unknown prices remain missing; public tasks do not imply free execution.
A useful planning expression is \[
 C_{\rm total}=\sum_{\rm calls}
  (r_{\rm in}t_{\rm in}+r_{\rm out}t_{\rm out}+c_{\rm tools})
  +C_{\rm retrieval}+C_{\rm compute}+C_{\rm review}+C_{\rm rework},
\] where token rates and token counts use matching units and cache
distinctions are retained when applicable.

For illustration, 100 independent mechanisms in each of six disjoint
task families, two variants per mechanism and four live arms give 4,800
assigned arm-variant cases. Three attempts and up to three internal
model calls per attempt yield a conservative envelope of 43,200 calls.
If every call is capped at 2,000 input and 500 output tokens, the token
envelope is 86.4 million input and 21.6 million output tokens. At prices
expressed per million tokens, the model-only component is
\(86.4r_{\rm in}+21.6r_{\rm out}\), before the other costs above. These
are deliberately explicit caps, not a spend forecast or an approved
budget \citep{exec13_live_planning}. A single-family pilot and fewer
justified comparisons can be much cheaper; reducing the study should
also reduce its claims.

Before any final model calls, the design must bind exact backend
revisions, generation settings, permitted endpoints, prompts, retrieval
corpus, numerical algorithms, final mechanism distribution and draw,
reference adjudication, primary contrasts, multiplicity allocation,
numerical thresholds and a priced resource ceiling. A fixed stopping
rule retains all assigned tasks. A global outage or exhausted budget is
recorded as an interruption with unexecuted cases visible, not repaired
by silently replacing difficult cases. A materially changed model or
prompt requires a new evaluation identity.

The recommendation is therefore to fund a bounded prototype comparison
only when its development evidence identifies a consequential
uncertainty and its expected information can change the next investment
choice. The proposal can already specify the experiment, expose its
likely precision and resource trade-offs, and reject unsupported
universal performance claims. Authentic bank experiments become
worthwhile when the resulting prototype leaves a specific transport,
assignment or operating question that public and generated evidence
cannot answer.

\label{exec13:live-design-end}
